\FloatBarrier
\clearpage
\section{结论}
\label{sec:conclusion}
研究中的下一步，常常由刚得到的证据决定。EAR 把这种推进方式落实到三个工作流中：问题发现根据近邻文献重新取舍候选；数学内循环沿评阅指出的缺口修订研究，外循环用完整研究结果检验执行策略；评审回应围绕疑虑寻找证据，保留有效草稿并继续补齐缺口。研究材料、判断依据和版本变化随过程留下，为继续工作提供基础。

历史记录给出了这些机制的初步证据。固定稿件子集在修订后获得更好的模型评阅结果，G3 所选策略在历史批次中取得更多通过结果和更低的每通过研究调用；回应评分预测则同时显示了二类区分能力与下降类识别的不足。每项结果都对应具体任务、分母和评价条件。采用已有策略与继续搜索策略的投入，也需要分别核算。

\begin{insight}
\textbf{让研究持续推进，把关键判断留给人}\par
研究者给方向、定标准、检查关键结果。EAR 在这些约束下继续检索、尝试和修订，交回新的产物，以及支持下一步判断的依据。
\end{insight}

最终要检验的，是完成可信研究所需的人时。后续评估需要在匹配任务和一致质量检查下，记录阅读、编排、检查、修复、失败返工与最终审阅的投入，并同时核算完整机器费用和日历周期。这样才能回答 EAR 的设计目标在实际使用中实现了多少。

\Sub{开源项目与公开材料}
\href{https://github.com/liyingze07-svg/EffectiveAutoResearch-EAR-}{EAR 开源仓库}提供问题发现、自主数学研究与评审回应三个模块。随报告公开的聚合数据、评价协议和矢量图支持核对统计与检查论述；材料的范围与限制见第\ref{sec:rebuttal-results} 节。读者可以据此试用对应流程，核对报告结果，并参与后续改进。\src{1,10}
